% from ega4-15-fr.tex
\oldpage[IV-3]{240}
noethérien), comme $X_y^0$ est la seule composante irréductible de $f^{-1}(y)$ contenant $x$, on déduit de (15.6.3) et (15.2.3) que si $\dim(X_y^0)$ est constante dans un voisinage de $y_0$, $f$ est plat au point $x_0$.

Pour prouver (iii), remarquons d'abord que l'ensemble $X^0$ est localement constructible dans $X$ (9.7.12), donc, par (9.5.5), l'ensemble $E$ des $y\in Y$ tels que $\dim(X_y^0)=\dim(X_{y_0}^0)$ est localement constructible dans $Y$. Appliquons alors (1.10.1) à $E$~: il suffit de prouver (compte tenu de (i)) que $\dim(X_{y'}^0)=\dim(X_{y_0}^0)$ pour le point générique $y'$ d'une composante irréductible de $Y$ contenant $y_0$~; ceci permet déjà de supposer que $Y=\operatorname{Spec}(\mathcal O_{Y,y_0})$.

Si on est dans le cas $\alpha$), on se ramène aussitôt, en vertu de (\textbf{II}, 7.1.9), et utilisant (4.5.13) et (4.4.1) comme dans (15.6.4), au cas où $Y$ est spectre d'un anneau de valuation discrète, $y_0$ son point fermé et $y'$ son point générique. Mais alors les hypothèses entraînent en vertu de (15.2.2.1), que $f$ est \emph{plat} au point $g(y_0)$, donc aussi dans un voisinage $V$ de ce point dans $X$ (11.1.1). Pour démontrer notre assertion, on peut remplacer $X$ par $V$, car $V\cap X_{y'}^0$ n'est pas vide en vertu de (2.3.4), donc est un ouvert partout dense dans $X_{y'}^0$ et a par suite même dimension (4.1.1.3)~; d'ailleurs, comme $X_{y'}^0$ est aussi la composante connexe de $f^{-1}(y')$ contenant $g(y')$, on peut supposer $V$ pris tel que $V$ ne rencontre aucune autre composante irréductible de $f^{-1}(y_0)$ ni de $f^{-1}(y')$, autrement dit on peut supposer que $X^0=X$~; mais alors la conclusion résulte de (12.1.1, (i)), puisque par hypothèse $X_{y_0}^0$ est intègre.

Supposons maintenant qu'on soit dans le cas $\beta$). Appliquons cette fois (\textbf{II}, 7.1.4) de la même manière que (\textbf{II}, 7.1.9) dans le cas $\alpha$)~: on est alors ramené au cas où $Y$ est spectre d'un anneau de valuation (non nécessairement discrète), $y_0$ son point fermé et $y'$ son point générique. En vertu de (14.3.13), il existe une composante irréductible $Z$ de $X$ contenant $X_{y_0}^0$ et dominant $Y$, et telle que $\dim(Z\cap f^{-1}(y'))=\dim(X_{y_0}^0)$~; mais l'hypothèse $\beta$) et le fait que $\dim(Z\cap f^{-1}(y'))\leqslant\dim(f^{-1}(y'))$ montrent que $\dim(X_{y_0}^0)\leqslant\dim(X_{y'}^0)$, d'où $\dim(X_{y_0}^0)=\dim(X_{y'}^0)$ en vertu de (i).

Reste à prouver (iv). Notons d'abord que les ensembles envisagés ne changent pas lorsqu'on remplace $f$ par $f_{(Y_{\mathrm{red}})}:X\times_Y Y_{\mathrm{red}}\to Y_{\mathrm{red}}$, la projection $X\times_Y Y_{\mathrm{red}}\to X$ étant un homéomorphisme. Autrement dit, on peut supposer $Y$ réduit, et alors il résulte de (ii) que $f$ est \emph{plat} aux points de $W$. Considérons un point $x_0$ de $W$ et prouvons que $W$ est un voisinage de $x_0$~; procédant comme au début de la démonstration, et utilisant aussi (11.2.7), on peut de plus supposer $Y$ noethérien~; il résulte alors de (ii) et de (15.6.4) que $X^0$ est un voisinage de $x_0$ dans $X$, et de (12.1.1, (vii)) que $W$ est aussi un voisinage de $x_0$ dans $X$.
\end{proof}

\begin{corollary}[15.6.7]
\label{IV.15.6.7-fr}
Supposons vérifiées les conditions préliminaires de (15.6.6) sur $f$, et supposons en outre que pour tout $y\in Y$, $X_y^0$ soit géométriquement intègre sur $\mathbf{k}(y)$. Alors les conditions suivantes sont équivalentes :
\begin{enumerate}[label=\emph{\alph*)}]
  \item La fonction $y\mapsto\dim(X_y^0)$ est localement constante dans $Y$.
  \item Le morphisme $f_{(Y_{\mathrm{red}})}:X\times_Y Y_{\mathrm{red}}\to Y_{\mathrm{red}}$ déduit de $f$ par changement de base, est plat aux points de $X^0$.
\end{enumerate}

